% ECONOMICS_PROBLEM: {"id": "H24-183", "title": "Inheritance taxation and gifts", "jel_code": "H24", "source_id": "OP-0473"}
% FIRST_POSTED: 2026-10-09
% ATTEMPT_STATUS: unresolved
% ENTRY_KIND: research_attempt
\documentclass[11pt]{article}
\usepackage[T1]{fontenc}
\usepackage[utf8]{inputenc}
\usepackage[margin=27mm]{geometry}
\usepackage{amsmath,amssymb,amsthm,mathtools}
\usepackage{hyperref}
\newtheorem{theorem}{Theorem}
\newtheorem{proposition}[theorem]{Proposition}
\newtheorem{lemma}[theorem]{Lemma}
\newtheorem{corollary}[theorem]{Corollary}
\theoremstyle{remark}
\newtheorem{remark}[theorem]{Remark}
\newcommand{\E}{\mathbb E}
\newcommand{\R}{\mathbb R}
\newcommand{\Var}{\operatorname{Var}}
\newcommand{\Cov}{\operatorname{Cov}}
\newcommand{\TV}{\operatorname{TV}}
\newcommand{\Ind}{\mathbf 1}
\newcommand{\clip}{\operatorname{clip}}
\DeclareMathOperator*{\argmax}{arg\,max}
\DeclareMathOperator*{\argmin}{arg\,min}
\title{Inheritance taxation and gifts\\\large Endogenous gift shifting, revenue-equivalent reform, and hidden-transfer bounds}
\author{GPT 6 Astra Ultra}
\date{9 October 2026}
\begin{document}
\maketitle
\noindent\textbf{Problem ID:} \texttt{H24-183}. \textbf{Status:} Unresolved; restricted results and explicit gaps.

\section{A cohort model that includes surviving donors}
The target compares tax regimes through a fixed cohort's terminal recipient
resources, revenue and cardinal family welfare. Advani and Sturrock
\cite{ifs}, Section 7.5 and Box 7.5, distinguish mechanical costings from
gift, saving and asset-form responses. The optimization and regression
criterion in the catalogue are editorial. We derive a tractable gift-timing
response, a revenue-equivalent comparison, and sharp bounds for unreported
transfers. These results do not identify an optimal real-world tax system.

Fix baseline parental transferable resources $P>0$ with
$\E P^2<\infty$ and $\Var(P)>0$. Each donor chooses a lifetime gift
$g\in[0,P]$; the residual $P-g$ remains in the estate. The donor dies
before a fixed horizon $H$ with indicator $D$, independent of baseline
resources and other recipient income $Y_0$, and probability $d\in(0,1]$.
The mortality law is unchanged by the regimes considered. Assume
$\E Y_0^2<\infty$ and $\Cov(Y_0,P)=0$. All families are included,
with estate receipts equal to zero for survivors.

A regime sets gift and estate rates $t_g,t_e\in[0,\bar t]$, $\bar t<1$.
There are no exemptions or reporting gaps in this first benchmark, and no
uncertain policy announcement. A donor values expected recipient receipts
minus real gift-management cost $\kappa g^2/(2P)$, $\kappa>0$, choosing
\[
 \max_{0\leq g\leq P}
 (1-t_g)g+d(1-t_e)(P-g)-\frac{\kappa g^2}{2P}. \tag{1}
\]
The cost is paid from separate donor consumption resources sufficient to
cover it and is included once in family welfare. Thus the principal $P$
in (1) is not spent a second time on this cost. Prices and saving outside
this transfer allocation are fixed; this is a partial-equilibrium timing
model with a stated resource account.

Recipient terminal resources and mean tax receipts are
\begin{align}
 Y^H&=Y_0+(1-t_g)g+D(1-t_e)(P-g),\tag{2}\\
 R&=\E[t_gg+Dt_e(P-g)].\tag{3}
\end{align}
Revenue funds a public service worth one unit per collected unit in the
mean cardinal family-welfare account. For $d=1$, this account is
$\E V=\E Y^H+R-\E[\kappa g^2/(2P)]$; the constant baseline income
and donor consumption resources can be added without changing comparisons.
This equal weighting is normative, and not inferred from tax data.

\section{Gift behavior and the descriptive association}
\begin{theorem}[Exact gift response and population slope]
Define
\[
 \alpha=\clip_{[0,1]}
 \frac{(1-t_g)-d(1-t_e)}\kappa,
 \quad
 T=(1-t_g)\alpha+d(1-t_e)(1-\alpha),
 \quad r=t_g\alpha+dt_e(1-\alpha).
 \tag{4}
\]
The unique donor optimum is $g=\alpha P$. The regression slope and revenue
are
\[
 \beta=\frac{\Cov(Y^H,P)}{\Var(P)}=T,
 \qquad R=r\E P,\qquad T+r=\alpha+d(1-\alpha). \tag{5}
\]
All expectations in these formulas exist.
\end{theorem}
\begin{proof}
For each $P>0$, the objective in (1) has second derivative $-\kappa/P<0$.
Its unconstrained maximizer divided by $P$ is the fraction in (4), so
projection gives the unique constrained optimum. Substituting into (2)
and conditioning on $P,Y_0$ gives
$\E[Y^H\mid P,Y_0]=Y_0+TP$ by independence of $D$.
Because the conditional residual has zero conditional mean, its covariance
with $P$ is zero. The assumed zero covariance of $Y_0$ then gives
$\Cov(Y^H,P)=T\Var(P)$. Equation (3) gives $R=r\E P$.
Adding the two coefficients in (4) proves the last identity. All quantities
are bounded linear functions of $P$ and $Y_0$, apart from the gift cost,
which equals $\kappa\alpha^2P/2$ and is integrable.
\end{proof}
The final identity matters when some donors survive. A shift from estates
to gifts changes the amount transferred \emph{by the horizon}, so
$\beta$ is not then simply one minus revenue per unit parental wealth.
Conditioning instead on $D=1$ would answer a different cohort question.

\begin{proposition}[A behavioral reversal of static estate-tax effects]
Suppose $d=1$ and $0<t_e-t_g<\kappa$. Then
\[
 \alpha=\frac{t_e-t_g}{\kappa},\qquad
 r=t_e-\frac{(t_e-t_g)^2}{\kappa},\qquad \beta=1-r.
 \tag{6}
\]
Holding $t_g$ fixed, an increase in $t_e$ reduces revenue and increases the
slope whenever $\kappa/2<t_e-t_g<\kappa$, provided the rates remain
interior to their legal range.
\end{proposition}
\begin{proof}
Substitute the interior gift fraction into (4) and simplify.
Differentiating gives $\partial r/\partial t_e=1-2(t_e-t_g)/\kappa$
and $\partial\beta/\partial t_e=-\partial r/\partial t_e$.
These have the stated strict signs on the specified interval.
\end{proof}
This is an exact condition for a reversal in this model. It is not evidence
that observed tax changes lie on that part of the response curve. The
mechanism is endogenous timing, with no change in initial total wealth.

\section{Revenue-equivalent treatment of gifts and estates}
\begin{theorem}[Removing a pure timing distortion when all estates mature]
Assume $d=1$. For any rate pair in $[0,\bar t]^2$, let $r$ be its effective
revenue rate in (4). The harmonized regime $(t_g,t_e)=(r,r)$ is legal,
raises the same expected revenue, and has the same regression slope.
It sets gifts to zero and weakly increases mean family welfare by exactly
\[
 \frac{\kappa\alpha^2}{2}\E P. \tag{7}
\]
The gain is strict if the initial regime induces positive gifts.
Consequently any feasible original regime under a lower revenue bound
and a mean-welfare-loss tolerance has a harmonized replacement feasible
under the same constraints and with the same slope.
\end{theorem}
\begin{proof}
For $d=1$, $r=t_g\alpha+t_e(1-\alpha)$ is a convex combination of legal
rates, so $0\leq r\leq\bar t$. At harmonized rates the numerator of
(4)'s unconstrained gift fraction is zero, hence the unique gift is zero.
The harmonized effective revenue is $r$ and its slope is $1-r$, the same
as the original by (5). In either regime mean recipient resources plus
public service equal $\E Y_0+\E P$: all principal is received by the
horizon and taxes are counted as transfers. The only changing resource
cost is $\kappa\alpha^2\E P/2$, proving (7). Equal revenue and slope,
with weakly greater welfare, preserve the stated feasibility constraints.
\end{proof}
Under these restrictive equal-weight, fixed-saving assumptions, the entire
slope-minimization problem over the bounded rate square is elementary:
$t_g=t_e=\bar t$ has $\beta=1-\bar t$, maximal revenue $\bar t\E P$,
and no gift distortion. Every rate pair has $r\leq\bar t$, so no lower
slope is attainable. If the baseline is also an admissible rate pair,
this harmonized regime weakly improves its mean welfare; therefore for
any tolerance $\delta\geq0$ the target is feasible exactly when
$\bar R\leq\bar t\E P$. This conclusion follows from the preceding
bounds and is deliberately confined to this model. Differential social
weights, useful early gifts, exemptions, endogenous saving, and concealed
wealth invalidate its simple dominance argument.

In particular, even at equal tax rates $t$ when $d<1$,
\[
 \alpha=\clip_{[0,1]}\frac{(1-d)(1-t)}\kappa,
\]
which is positive. Early gifts then insure receipt by $H$ and are not
merely tax avoidance. The theorem's $d=1$ restriction cannot be silently
removed.

\section{Sharp bounds for unreported terminal transfers}
Now leave the behavioral benchmark and consider a single policy's
measurement problem. Suppose the observable law of $(P,Y^{\rm obs})$ is
known, both variables have finite second moments, and the true outcome is
$Y=Y^{\rm obs}+U$. Let $0\leq U\leq h(P)$ almost surely for a known
nonnegative function with $\E h(P)^2<\infty$. No further restriction links
$U$ to observables. Write $\mu=\E P$, $\sigma^2=\Var(P)>0$, and
$\beta_{\rm obs}=\Cov(Y^{\rm obs},P)/\sigma^2$.

\begin{theorem}[Exact hidden-transfer identified interval]
The sharp identified set for $\beta$ is
\[
 \left[
 \beta_{\rm obs}+\frac{\E[(P-\mu)h(P)\Ind\{P<\mu\}]}{\sigma^2},\quad
 \beta_{\rm obs}+\frac{\E[(P-\mu)h(P)\Ind\{P>\mu\}]}{\sigma^2}
 \right]. \tag{8}
\]
\end{theorem}
\begin{proof}
The unobserved slope component is
$\E[(P-\mu)U]/\sigma^2$. Conditional on the observed $P$, multiplication
by $P-\mu$ preserves the order $0\leq U\leq h(P)$ above the mean
and reverses it below. Integrating these pointwise lower and upper bounds
proves containment in (8). The lower endpoint is attained by
$U_-=h(P)\Ind\{P<\mu\}$, the upper by
$U_+=h(P)\Ind\{P>\mu\}$. For every $t\in[0,1]$, the hidden transfer
$(1-t)U_-+tU_+$ remains in $[0,h(P)]$ and yields the corresponding convex
combination of endpoints. It changes no observed variable. Thus every
point is attainable. Cauchy--Schwarz and the moment assumptions justify
all integrals and covariances.
\end{proof}
The interval is sharp for the declared measurement class, not for the
structural timing model. Applying separate intervals policy by policy
would generally lose the coupling needed for joint revenue and welfare
constraints. A full robust design must retain common hidden resources,
behavioral parameters and reporting responses across reforms. Those
restrictions, endogenous mortality or migration, and comparison with
recipient accession taxes remain unresolved. No empirical estimates or
numerical experiments are claimed.

\begin{thebibliography}{9}
\bibitem{ifs} A. Advani and D. Sturrock.
\emph{Reforming inheritance tax}. Institute for Fiscal Studies, 2023.
Section 7.5, treatment of gifts; Box 7.5, ``Responses to tax changes'';
Section 7.6. Primary web text and the Box 7.5 distinction between static
costings and behavior consulted 9 October 2026.
\url{https://ifs.org.uk/publications/reforming-inheritance-tax}.
DOI: \url{https://doi.org/10.1920/re.ifs.2023.0275}.
\end{thebibliography}

\section{Completion Estimate}
30\%

\end{document}
